\section{Funções contínuas}

\begin{definition}
    Uma função \(f: X \to Y \) é contínua num ponto \(a \in X\) se \(\forall V\ab \subset Y : f(a) \in V\), \(\exists U\ab\subset X : a \in U\) e \(f(U)\subset V\).
\end{definition}

\begin{note}
    \(f\) é contínua em geral se for contínua em cada ponto. Denotamos \(C(X,Y)\) o conjunto de todas as funções contínuas  e \(C(X)\) se \(Y = \R\).
\end{note}

\begin{proposition}
    \textbf{8.3.} Dada \(f : X \to Y\) são equivalentes:
    \vspace{-0.3cm}
    \begin{enumerate}[label =(\alph*), left = 0.6cm]
        \item \(f\in C(X,Y)\). \hspace{3.9cm} (b) \ \(\forall V \in \U_{f(a)}, \ \exists U \in \U_a : f(U) \subset V\). 
        \item[(c)] \(f^{-1}(V)\ab \subset X, \ \forall V\ab\subset Y \). \hspace{1.3cm} (d) \ \(f^{-1}(B)\ce \subset X, \ \forall B\ce\subset Y \).
    \end{enumerate}
\end{proposition}

\begin{proposition}
    \textbf{8.4.} Seja \(g \circ f: X \overset{f}{\longrightarrow} Y \overset{g}{\longrightarrow} Z\). Se \(f\) é contínua em \(a\) e \(g\) é contínua em \(f(a)\), então \(g \circ f\) é contínua em \(a\). 
\end{proposition}

\begin{proposition}
    \textbf{8.6.} Se \(f \in C(X,Y)\) e \(S \leq X\), então \(f|_S \in C(S, Y)\). 
\end{proposition}

\begin{proposition}
    \textbf{8.7.} Sejam \(X = S_1 \cup S_2\), ambos abertos (ou fechados) e \(f:X \to Y\) tal que \( f|_{S_j}\in C(S_j,Y)\), então \(f \in C(X,Y)\). 
\end{proposition}

\demo{
    \-\hfill  \(f^{-1}(V\ab) = (S_1 \cap f^{-1}(V)) \cup (S_2 \cap f^{-1}(V)) = (S_1 \cap U_1\ab) \cup (S_2\cap U_2\ab)\) \-\hfill
}

\begin{definition}
    \(f: X \to Y\) é um \emph{homeomorfismo} se \(f\) for bijetora, contínua e \(f^{-1}\) tambem for contínua, denotamos \(f: X \simeq Y\); \emph{mergulho} se \(f: X \simeq \operatorname{Im}f\leq Y\).    
\end{definition}

\begin{definition}
    \(f\) é \emph{aberta} se \(f(U)\ab\subset Y, \ \forall U\ab\subset X\); \emph{fechada} se \(f(B)\ce\subset Y, \ \forall B\ce\subset X\). 
\end{definition}

\begin{proposition}
    \textbf{8.9.} Dada \(f: X \to Y\) bijetiva, são equivalentes: 
    \vspace{-0.3cm}
    \begin{enumerate}[label = (\alph*), left = 0.6cm]
        \item \(f:X \simeq Y\). \hfill (b) \ \(f\) é contínua e aberta. \hfill (c) \ \(f\) é contínua e fechada. 
    \end{enumerate}
\end{proposition}

\subsection*{Produtos infinitos e axioma de escolha}

\begin{definition}
    Sejam \((X_i)_{i\in I}\) familia não vazia de conjuntos, o seu \emph{produto cartesiano} é 
    \[\prod_{i\in I}X_i := \left\{\x : I \to \bigcup_{i\in I} X_i \mid \forall i \in I, \ \x(i) = x_i \in X_i\right\}.\]
    Podemos escrever a conveniência \(\x = (x_i)_{i\in I} \in \prod X_i\). Dado \(j \in I\) a \emph{projeção} \(j\)\emph{-ésima} é a função tal que \(\pi_j : \prod X_i \ni (x_i) \mapsto x_j \in X_j\). 
\end{definition}

\begin{note}
    Mesmo que \(\forall i \in I, \ X_i\neq \varnothing\) não tem como garantir que \(\prod X_i \neq \varnothing\).
\end{note}

\begin{definition}[\emph{Axioma da escolha}]
    Sejam \((X_i)_{i\in I}\) familia não vazia de conjuntos disjuntos não vazios. Então, \(\exists f: I \to \bigcup X_i\) tal que \(\forall i \in I, \ f(i) \in X_i\) (\emph{função de escolha}).
\end{definition}

\begin{proposition}
    \textbf{9.3.} Se \((X_i)_{i \in I}\) é familia não vazia de conjuntos não vazios, então \(\prod X_i \neq \varnothing\). \comentario{Se os \(X_i\) não fossem disjuntos, considera \(Y_i := X_i \times \{i\}\) e usa o axioma } 
\end{proposition}

\subsection*{O espaço produto}

\begin{note}
    Daquí pra frente, \((X_i)_{i\in I}\) é familia não vazia de conjuntos não vazios. 
\end{note}

\begin{proposition}
    \textbf{10.1.} A familia \(\B := \left\{\prod U_i : \forall i \in I, \  U_i\ab \subset X_i\right\}\) é base pra uma certa topologia em \(\prod X_i\), que chamamos de \emph{topologia das caixas}.\footnote{Mesmo sendo fácil de definir, essa topología não é muito conveniente.} \comentario{Usa \(\square \) \hyperlink{prop66}{6.6.}}
\end{proposition}

\begin{proposition}
    \textbf{10.2.} Seja \(\B\) a familia de todos os \(\prod U_i  \) tais que:
    \vspace{-0.3cm}
    \begin{itemize}[left = 1.8cm]
        \item[(TP1)] \ \(\forall i \in I, \ U_i\ab \subset X_i\). \hfill  (TP2) \ \( \forall i \in I\setminus J, \ U_i = X_i,  \ J\subset I \text{ finito}\). 
    \end{itemize}
    \vspace{-0.3cm}
    Então, \(\B\) é base de uma certa topologia em \(\prod X_i\), que chamamos de \emph{topologia produto}.
\end{proposition}

\demo{
    Mais uma vez usa \(\square\) \hyperlink{prop66}{6.6.}, é conveniente escrever \(U\ab \in \B \) como sendo, 
    \[U = \prod_{j \in J} U_j\ab \times \prod_{i \in I\setminus J} X_i = \bigcap_{j\in J} \pi^{-1}_j(U_j)\ab.\]
    \-\vspace{-0.3cm}
}

\begin{note}
    A menos de aclaração, consideramos sempre \( \prod X_i\) com a topologia produto. 
\end{note}
\begin{proposition}
    \textbf{10.3.} A topologia producto é a topologia mais fraca que faz das projeções \(\pi_j : \prod X_i \to X_j\) funções contínuas. \comentario{Suponha outra \(\tau\) e mostra que \(\tau_{\operatorname{prod}} \subset \tau\)}
\end{proposition}

\begin{proposition}
    \textbf{10.4.}  \(g : Y \to \prod X_i\) é contínua \(\leftrightharpoons \ \forall j \in I, \ \pi_j\circ g : Y \to X_j\) é contínua. \comentario{Pra a volta trabalha a expressão como interseções} 
\end{proposition}

\begin{proposition}
    \textbf{10.5.} Sejam \(X\) um conjunto e \((f_i : X \to X_i)_{i\in I}\) uma familia de funções e 
    \[\B := \left\{\bigcap_{j\in J} f_j^{-1}(U_j) : J \subset I \ \text{ é finito e } \forall j \in J, \  U_j\ab \subset X_j\right\}.\] 
    \vspace{-0.8cm}
    \begin{enumerate}[label = (\alph*), left = 0.6cm]
        \item \( \B\) é base pra uma topologia \(\tau_w\) em \(X\), que chamamos de \emph{topologia fraca}.
        \item \(\tau_w\) (\emph{weak topology}) é a topologia mais fraca que faz as \((f_i)\) contínuas.
        \item \(g:Y \to X\) é contínua \(\leftrightharpoons \ \forall j \in I, \ f_j \circ g\) é contínua. 
    \end{enumerate}
\end{proposition}

\subsection*{O espaço quociente}

\begin{proposition}
    \textbf{11.1.} Sejam \((X, \tau)\), \(Y\) um conjunto qualquer e \(\pi : X \twoheadrightarrow Y\). Então, a coleção 
    \[\tau_\pi := \{V\subset Y : \pi^{-1}(V)\ab \subset X\}\] 
    é uma topologia em \(Y\) que chamamos de \emph{topologia quociente} dada por \(\pi\).
\end{proposition}

\begin{definition}
    Uma aplicação \(\pi: (X,\tau_X) \twoheadrightarrow (Y, \tau_Y)\) é \emph{quociente} se \(\tau_Y = \tau_\pi\).  
\end{definition}

\begin{proposition}
    \textbf{11.3.} A topologia \(\tau_\pi\) é a mais fraca tal que \(\pi: X \to (Y,\tau_\pi)\) é contínua. 
\end{proposition}

\begin{proposition}
    \textbf{11.4.} Sejam \(\pi : X \to Y\) aplicação quociente e \(g: Y \to (Z,\tau_Z)\) uma função. Então, \(g\) é contínua \ \(\leftrightharpoons\) \ \(g\circ \pi: X \to Z\) é contínua.
\end{proposition}

\begin{proposition}
    \textbf{11.5.} Seja \(\pi: X \twoheadrightarrow Y\) contínua. Se \(\pi\) é aberta (ou fechada), então \(\tau_Y = \tau_\pi\). \comentario{É coisa de provar as duas contenções, não é díficil}
\end{proposition}

\begin{example}
    Seja \(\pi : [0,2\pi] \to \E^1\) tal que  \(t \mapsto (\cos t, \sin t)\). \comentario{Sobrejetiva e fechada}
\end{example}

\begin{proposition}
    \textbf{11.7.} Sejam \((X, \tau_X)\), \ \( \mathcal{D} \subset \wp(X) : X = \bigsqcup \{A : A \in \mathcal{D}\}\) (\emph{decomposição}) e  
    \[\tau_{\mathcal{D}} = \left\{\mathcal{A}\subset \mathcal{D} : \left(\bigcup \{A : A \in \mathcal{A}\}\right)\ab \subset X\right\}.\]
    Então \(\tau_{\mathcal{D}}\) é uma topologia em \(\mathcal{D}\). \comentario{Não é díficil}
\end{proposition}

\begin{note}
    Dado \(x \in X, \ \exists ! A_x \in \mathcal{D} : x \in A_x \), logo \(P : X \to \mathcal{D}\) tal que \(x \mapsto A_x\) é bem definida, chamamos ela de \emph{aplicação decomposição}. 
\end{note}

\begin{proposition}
    \textbf{11.8.} Toda aplicação decomposição é uma aplicação quociente.  
\end{proposition}

\begin{proposition}
    \textbf{11.9.} Dada \(\pi\) aplicação quociente, existem \(P\) aplicação decomposição e \(f: Y \simeq \mathcal{D}  \) tal que \(f \circ \pi = P\). \comentario{Olha \cite[pag. 30]{mujica2005}}
\end{proposition}

\begin{definition}
    Dados \(X\) e \(\sim \) relação de equivalência em \(X\), o espaço \(\nicefrac{X}{\sim}\) é chamado de \emph{espaço de identificação} de \(X\) módulo \(\sim\). 
\end{definition}



